\chapter*{Preface}
\addcontentsline{toc}{chapter}{Preface}
\markboth{Preface}{Preface}

The first edition of this book grew out of one simple problem: the collection of
textbooks I had gathered ran to more than ten thousand pages, but those pages were
never read in sequence. Opening Rudin at chapter three without ever touching the
definition of a supremum produces nothing but frustration. What was missing was not
material, but order.

Before this manuscript there was a thirteen-chapter version. Read again, it showed two
weaknesses. First, thirteen chapters went too fast, so the hardest parts went by in a
glance. Second, that book could be read by only one kind of reader: a student already
used to formal notation. A curious middle-school student had no door in, and a lecturer
found no teaching notes inside it.

This version repairs both. Its scope has been widened to six parts and thirty-two
chapters, with material the thirteen-chapter version did not have: sets and
cardinality, infinite series, systems of linear equations and determinants, multiple
integrals together with the vector field theorems, differential equations and dynamical
systems, regression, and model generalization and regularization. Every chapter is now
built in layers.

The first layer is the \emph{For younger readers} box. It holds analogies to everyday
objects, without a single mathematical symbol, and closes with one question that makes
the reader curious. The second layer is the body of the chapter itself: definitions,
theorems, proofs, worked examples, and exercises, as in an ordinary textbook. The third
layer is the \emph{Note for instructors} box: the teaching order I recommend, the
misconceptions that almost always surface, and one classroom activity. The fourth layer
is the \emph{For researchers} box: what is really proved, what is only popularized, the
primary sources, and the questions that are still open.

The four layers do not replace one another: a younger reader can read a chapter from the
orange box and the summary alone, an undergraduate reads everything except the
researcher box, an instructor uses the green box as a plan for the meeting, and a
researcher starts from the gray box. One book, four paths.

Those layers set the \emph{depth}. What sets the \emph{way in} is the second set of four
boxes: \emph{Read it aloud} for readers who understand by hearing, \emph{Do it by hand}
for those who understand by moving and touching, \emph{Step-by-step recipe} for those who
understand by memorizing a procedure first, and \emph{Challenge} for those who understand
precisely when the problem is not yet clear. The two sets live together: one and the same
chapter, with six doors.

Three principles of writing have not changed. \emph{No claim without a testing tool}:
every theorem comes with a proof, or with an honest admission that the proof is skipped
together with where it is skipped; every number comes from a real execution that the
reader can repeat. \emph{Notation is introduced before it is used}. \emph{Honesty in
evaluation is part of mathematics}: the last chapters do not stop at how to build a
model, but at how to know whether a claim about that model is proved.

One thing is deliberately absent from this book: complete answer keys. Giving a final
number without a trace of the check would violate the first principle. In its place,
Appendix~C gives grading criteria and a verification tool for every kind of exercise, so
that readers can check their own work in the same way the author checked this manuscript.

This book is not a substitute for the source texts it points to, but an ordered way in:
if one chapter makes a reader open Rudin, Axler, Boyd, or a paper by Athalye and then
read the section it names, this book has succeeded. Whatever errors remain belong to the
author, not to those texts.

\enlargethispage{3\baselineskip}
\vspace{0.35cm}
\begin{flushright}
Afiq Andico Pangimpian\\
October 2026
\end{flushright}
